% =========================================================
% Topic 3.7
% ডিজিটাল ডিভাইসের পরিচিতি
% Introduction to Digital Device
% =========================================================

\section{ডিজিটাল ডিভাইসের পরিচিতি
(\textenglish{Introduction to Digital Device})}

বর্তমান যুগে আমাদের দৈনন্দিন জীবন, শিক্ষা, চিকিৎসা, ব্যাংকিং,
যোগাযোগ, শিল্পকারখানা ও বিনোদনের প্রায় প্রতিটি ক্ষেত্রেই
ডিজিটাল প্রযুক্তির ব্যবহার রয়েছে।

মোবাইল ফোন, কম্পিউটার, ল্যাপটপ, ডিজিটাল ঘড়ি,
\textenglish{ATM}, ডিজিটাল ক্যামেরা, প্রিন্টার,
টেলিভিশন, বিভিন্ন চিকিৎসা-যন্ত্র এবং স্বয়ংক্রিয় নিয়ন্ত্রণব্যবস্থা—
এসবই ডিজিটাল প্রযুক্তিনির্ভর যন্ত্রের উদাহরণ।

ডিজিটাল ব্যবস্থায় তথ্যকে সাধারণত বিচ্ছিন্ন অবস্থার মাধ্যমে প্রকাশ করা হয়।
কম্পিউটার ও অধিকাংশ ডিজিটাল ইলেকট্রনিক সার্কিটে
এই অবস্থাগুলোকে প্রধানত \textbf{০ এবং ১} দ্বারা প্রকাশ করা হয়।

সাধারণভাবে—

\[
0 \rightarrow \text{\textenglish{OFF/Low}}
\]

\[
1 \rightarrow \text{\textenglish{ON/High}}
\]

অর্থাৎ ডিজিটাল সার্কিটে দুটি পৃথক বৈদ্যুতিক অবস্থার মাধ্যমে
তথ্য প্রকাশ ও প্রক্রিয়াকরণ করা হয়।

% =========================================================
% What is Digital Device
% =========================================================

\subsection{ডিজিটাল ডিভাইস কী?
(\textenglish{What is a Digital Device?})}

\par\begingroup

\textbf{সংজ্ঞা:}

যে ইলেকট্রনিক যন্ত্র বা সার্কিট ডিজিটাল সংকেত ব্যবহার করে
তথ্য গ্রহণ, প্রক্রিয়াকরণ, সংরক্ষণ অথবা ফলাফল প্রদান করে,
তাকে \textbf{ডিজিটাল ডিভাইস
(\textenglish{Digital Device})} বলে।
\par\endgroup


\textbf{\textcolor{ICTOrange}{উদাহরণ:}}

মোবাইল ফোন, কম্পিউটার, ক্যালকুলেটর, ডিজিটাল ঘড়ি,
\textenglish{ATM}, ডিজিটাল ক্যামেরা ইত্যাদি।

% =========================================================
% Main Working Flow
% =========================================================

\subsection{ডিজিটাল ডিভাইসের কাজের মূল ধারা}

একটি ডিজিটাল ডিভাইস সাধারণভাবে তিনটি ধাপে কাজ করে—

\[
\boxed{
\text{\textenglish{Input}}
\rightarrow
\text{\textenglish{Processing}}
\rightarrow
\text{\textenglish{Output}}
}
\]

\subsubsection*{উদাহরণ: ক্যালকুলেটরে $7+4$}

ক্যালকুলেটরে $7+4$ ইনপুট দিলে ডিভাইসটি সংখ্যাগুলোকে
অভ্যন্তরীণভাবে বাইনারি সংকেতে প্রকাশ করে।

\[
7_{10}=0111_2
\]

\[
4_{10}=0100_2
\]

এখন যোগ করলে—

\[
0111_2+0100_2=1011_2
\]

এবং,

\[
1011_2=11_{10}
\]

তাই ক্যালকুলেটরের ডিসপ্লেতে ফলাফল হিসেবে
\textbf{11} দেখা যায়।

% =========================================================
% Why 0 and 1
% =========================================================

\subsection{কেন ০ ও ১ ব্যবহার করা হয়?}

ইলেকট্রনিক সার্কিটে দুটি পৃথক বৈদ্যুতিক অবস্থা
তৈরি ও শনাক্ত করা তুলনামূলক সহজ।

যেমন—
\begin{itemize}[label={{\bulletfont\textbullet}}, leftmargin=*, nosep]
    \item বিদ্যুৎ প্রবাহ নেই বা কম $\rightarrow 0$
    \item বিদ্যুৎ প্রবাহ আছে বা বেশি $\rightarrow 1$
\end{itemize}

এই কারণে বাইনারি পদ্ধতি ডিজিটাল ইলেকট্রনিক্সে
অত্যন্ত উপযোগী।

\par\begingroup

\textbf{মনে রাখবে:}

\[
\boxed{0,\;1}
\]

ডিজিটাল ইলেকট্রনিক্সে বহুল ব্যবহৃত দুটি যৌক্তিক অবস্থা।
\par\endgroup


% =========================================================
% Components
% =========================================================

\subsection{ডিজিটাল ডিভাইসের অভ্যন্তরে গুরুত্বপূর্ণ উপাদান}

একটি ডিজিটাল ডিভাইসের ভেতরে বিভিন্ন ইলেকট্রনিক উপাদান
ও সার্কিট থাকে। গুরুত্বপূর্ণ কয়েকটি হলো—

\begin{enumerate}
    \item \textenglish{Transistor}
    \item \textenglish{Logic Gate}
    \item \textenglish{Integrated Circuit (IC)}
    \item \textenglish{Encoder}
    \item \textenglish{Decoder}
    \item \textenglish{Adder}
    \item \textenglish{Register}
    \item \textenglish{Counter}
\end{enumerate}

এগুলো একসঙ্গে কাজ করে বিভিন্ন ধরনের ডিজিটাল ব্যবস্থা তৈরি করে।

% =========================================================
% Transistor
% =========================================================

\subsection{ট্রানজিস্টর
(\textenglish{Transistor})}

ডিজিটাল ইলেকট্রনিক্সের অন্যতম মৌলিক উপাদান হলো
\textenglish{Transistor}।

এটি বিভিন্ন কাজে ব্যবহৃত হলেও ডিজিটাল সার্কিটে
ট্রানজিস্টরকে অনেক সময় একটি ইলেকট্রনিক
\textenglish{Switch} হিসেবে ব্যবহার করা হয়।

ট্রানজিস্টরের দুটি অবস্থা ধরা যায়—

\[
\text{\textenglish{OFF}}\rightarrow0
\]

\[
\text{\textenglish{ON}}\rightarrow1
\]

অর্থাৎ ট্রানজিস্টরের \textenglish{ON/OFF} অবস্থার সাহায্যে
বাইনারি ০ ও ১ প্রকাশ করা যায়।

\par\begingroup

\textbf{সহজ ধারণা:}

অনেকগুলো \textenglish{Transistor} নির্দিষ্ট নিয়মে সংযুক্ত করে
\textenglish{Logic Gate} এবং আরও জটিল ডিজিটাল সার্কিট তৈরি করা যায়।
\par\endgroup


\[
\boxed{
\text{\textenglish{Transistor}}
\rightarrow
\text{\textenglish{Logic Gate}}
\rightarrow
\text{\textenglish{Digital Circuit}}
}
\]

% =========================================================
% Logic Gate
% =========================================================

\subsection{লজিক গেট
(\textenglish{Logic Gate})}

\par\begingroup

\textbf{সংজ্ঞা:}

যে ইলেকট্রনিক সার্কিট এক বা একাধিক বাইনারি ইনপুট গ্রহণ করে
নির্দিষ্ট যৌক্তিক নিয়ম অনুযায়ী একটি বাইনারি আউটপুট প্রদান করে,
তাকে \textbf{\textenglish{Logic Gate}} বলে।
\par\endgroup


\textenglish{Logic Gate}-কে ডিজিটাল সার্কিটের মৌলিক
\textenglish{Building Block} বলা হয়।

প্রধান \textenglish{Logic Gate} হলো—

\begin{itemize}[label={{\bulletfont\textbullet}}, leftmargin=*, nosep]
    \item \textenglish{AND}
    \item \textenglish{OR}
    \item \textenglish{NOT}
    \item \textenglish{NAND}
    \item \textenglish{NOR}
    \item \textenglish{XOR}
    \item \textenglish{XNOR}
\end{itemize}

\textbf{\textcolor{ICTOrange}{উদাহরণ:}}

\textenglish{AND Gate}-এ সব ইনপুট ১ হলে আউটপুট ১ হয়।

\[
1\cdot1=1
\]

অন্যথায় আউটপুট ০ হয়।

\par\begingroup

\textenglish{Logic Gate} নিয়ে পরবর্তী অংশে বিস্তারিত আলোচনা করা হবে।
\par\endgroup


% =========================================================
% IC
% =========================================================

\subsection{\textenglish{Integrated Circuit (IC)}}

\par\begingroup

\textbf{সংজ্ঞা:}

একটি ছোট \textenglish{Semiconductor Chip}-এর মধ্যে
বহু ইলেকট্রনিক উপাদান ও সার্কিট একত্রে তৈরি করা হলে
তাকে \textbf{\textenglish{Integrated Circuit (IC)}} বলে।
\par\endgroup


একটি \textenglish{IC}-এর মধ্যে অনেক
\textenglish{Transistor}, \textenglish{Logic Gate}
এবং অন্যান্য ইলেকট্রনিক উপাদান থাকতে পারে।

বর্তমানের \textenglish{Processor}, \textenglish{Memory},
\textenglish{Controller} এবং বিভিন্ন
\textenglish{Digital System} তৈরিতে
\textenglish{IC} অত্যন্ত গুরুত্বপূর্ণ।

\textbf{\textcolor{ICTOrange}{উদাহরণ:}}

কম্পিউটার, মোবাইল ফোন, ক্যালকুলেটর ও অন্যান্য
ডিজিটাল যন্ত্রে বিভিন্ন ধরনের \textenglish{IC} ব্যবহৃত হয়।

% =========================================================
% Relation
% =========================================================

\subsection{ট্রানজিস্টর, লজিক গেট ও \textenglish{IC}-এর সম্পর্ক}

ডিজিটাল ইলেকট্রনিক্সের গঠনকে সহজভাবে এভাবে দেখানো যায়—

\[
\boxed{
\text{\textenglish{Transistor}}
\rightarrow
\text{\textenglish{Logic Gate}}
\rightarrow
\text{\textenglish{IC/Digital Circuit}}
\rightarrow
\text{\textenglish{Digital Device}}
}
\]

অর্থাৎ \textenglish{Transistor} থেকে
\textenglish{Logic Gate},
\textenglish{Logic Gate} থেকে জটিল
\textenglish{Digital Circuit}
এবং এসব সার্কিট ব্যবহার করে বিভিন্ন
\textenglish{Digital Device} তৈরি করা হয়।

% =========================================================
% Everyday Use
% =========================================================

\subsection{দৈনন্দিন জীবনে ডিজিটাল ডিভাইসের ব্যবহার}

ডিজিটাল ডিভাইসের ব্যবহার বর্তমানে অত্যন্ত ব্যাপক।

যেমন—

\begin{itemize}[label={{\bulletfont\textbullet}}, leftmargin=*, nosep]
    \item মোবাইল ফোন
    \item কম্পিউটার ও ল্যাপটপ
    \item ডিজিটাল ঘড়ি
    \item \textenglish{ATM}
    \item ডিজিটাল ক্যামেরা
    \item প্রিন্টার
    \item স্মার্ট টেলিভিশন
    \item ক্যালকুলেটর
    \item বিভিন্ন চিকিৎসা-যন্ত্র
    \item শিল্পকারখানার স্বয়ংক্রিয় নিয়ন্ত্রণব্যবস্থা
\end{itemize}

% =========================================================
% Digital Device and Boolean Algebra
% =========================================================

\subsection{ডিজিটাল ডিভাইস ও
\textenglish{Boolean Algebra}}

ডিজিটাল সার্কিটে ০ ও ১-এর ওপর বিভিন্ন যৌক্তিক কাজ করা হয়।
এই যৌক্তিক কাজগুলো বিশ্লেষণ ও প্রকাশের জন্য
\textenglish{Boolean Algebra} ব্যবহার করা হয়।

অন্যদিকে \textenglish{Boolean Algebra}-এর বিভিন্ন যৌক্তিক
\textenglish{Operation}-কে বাস্তবে ইলেকট্রনিক সার্কিটের মাধ্যমে
বাস্তবায়ন করে \textenglish{Logic Gate}।

তাই \textenglish{Digital Device} বুঝতে হলে
পরবর্তী ধাপে \textenglish{Boolean Algebra} ও
\textenglish{Logic Gate} সম্পর্কে জানা প্রয়োজন।

\[
\boxed{
\text{\textenglish{Boolean Algebra}}
\rightarrow
\text{\textenglish{Logic Gate}}
\rightarrow
\text{\textenglish{Digital Circuit}}
}
\]

% =========================================================
% Summary
% =========================================================

\subsection{এক নজরে}

\begin{center}
\begin{tabularx}{\textwidth}{|L{0.28\textwidth}|Y|}
\hline
\rowcolor{ICTLightBlue}
\textbf{বিষয়} &
\textbf{মূল ধারণা} \\
\hline

\textenglish{Digital Device} &
ডিজিটাল সংকেত ব্যবহারকারী ইলেকট্রনিক যন্ত্র \\
\hline

\textenglish{Binary State} &
০ ও ১ \\
\hline

\textenglish{Transistor} &
\textenglish{Electronic Switch} হিসেবে কাজ করতে পারে \\
\hline

\textenglish{Logic Gate} &
বাইনারি ইনপুটের ওপর যৌক্তিক কাজ করে \\
\hline

\textenglish{IC} &
ছোট \textenglish{Chip}-এ বহু ইলেকট্রনিক উপাদান ও সার্কিট \\
\hline

মূল ধারা &
\textenglish{Input}
$\rightarrow$
\textenglish{Processing}
$\rightarrow$
\textenglish{Output} \\
\hline

\end{tabularx}
\end{center}

% =========================================================
% Final Remember Box
% =========================================================

\par\begingroup

\textbf{মনে রাখবে}

\[
\boxed{
\text{\textenglish{Input}}
\rightarrow
\text{\textenglish{Processing}}
\rightarrow
\text{\textenglish{Output}}
}
\]

\[
\boxed{
0\leftrightarrow\text{\textenglish{OFF}},
\qquad
1\leftrightarrow\text{\textenglish{ON}}
}
\]

\[
\boxed{
\text{\textenglish{Transistor}}
\rightarrow
\text{\textenglish{Logic Gate}}
\rightarrow
\text{\textenglish{Digital Circuit}}
}
\]
\par\endgroup


% =========================================================
% Exam-ready Definition
% =========================================================

\par\begingroup
\smallhead{\textbf{পরীক্ষার জন্য সংজ্ঞা}}


ডিজিটাল সংকেত ব্যবহার করে তথ্য গ্রহণ, প্রক্রিয়াকরণ,
সংরক্ষণ অথবা ফলাফল প্রদানকারী ইলেকট্রনিক যন্ত্র বা সার্কিটকে
\textbf{ডিজিটাল ডিভাইস
(\textenglish{Digital Device})} বলে।

\par\endgroup

\subsection{অ্যানালগ থেকে ডিজিটাল এবং সার্কিটের উপাদান}
অ্যানালগ সংকেতে মান ধারাবাহিকভাবে বদলাতে পারে; ডিজিটাল সংকেতে নির্দিষ্ট বিচ্ছিন্ন মান ব্যবহার হয়। বাস্তব শব্দ বা তাপমাত্রা নিজে বাইনারি বিট নয়। সেন্সর থেকে পাওয়া অ্যানালগ সংকেতকে প্রয়োজনমতো sampling ও মানের ধাপ নির্ধারণের পরে digital code-এ প্রকাশ করা হয়। এই রূপান্তরের সার্কিট হলো ADC; বিপরীত কাজের সার্কিট DAC।

\begin{center}\begin{tabularx}{\linewidth}{|L{.2\linewidth}|Y|Y|}\hline
বিষয় & অ্যানালগ & ডিজিটাল\\\hline
মান & নির্দিষ্ট সীমায় ধারাবাহিক & নির্ধারিত বিচ্ছিন্ন মান\\\hline
উদাহরণ & মাইক্রোফোনের voltage & digitized audio samples\\\hline
প্রক্রিয়াকরণ & অ্যানালগ সার্কিটে & logic gate/processor-এ\\\hline
সীমা & noise ও যন্ত্রের নির্ভুলতা & noise margin, sampling ও bit resolution\\\hline
\end{tabularx}\end{center}

\smallhead{ট্রানজিস্টর, switch ও IC}
\textbf{ট্রানজিস্টর} একটি semiconductor উপাদান; নিয়ন্ত্রণ সংকেত দিয়ে বিদ্যুৎ প্রবাহ চালু বা বন্ধ করার switch হিসেবে ব্যবহার করা যায়। কয়েকটি ট্রানজিস্টরের উপযুক্ত সংযোগে logic gate তৈরি হয়। বহু উপাদান ও সংযোগ একটি ছোট semiconductor chip-এ তৈরি হলে তাকে \textbf{সমন্বিত বর্তনী (Integrated Circuit, IC)} বলে। একটি IC-তে gate, flip-flop বা আরও জটিল circuit থাকতে পারে; IC একটি নির্দিষ্ট gate-এর নাম নয়।

Logic ০ মানে সবসময় শূন্য volt বা বিদ্যুৎ নেই, এবং logic ১ মানে সব যন্ত্রে একই নির্দিষ্ট voltage—এমন নয়। গেটের নির্ধারিত LOW/HIGH voltage range এবং polarity অনুযায়ী binary অর্থ নির্ধারিত হয়।

\chThreeExample{১}{ডিজিটাল থার্মোমিটারে তথ্যের পথ}
\chThreeAnswer সেন্সর তাপমাত্রা অনুযায়ী একটি বৈদ্যুতিক সংকেত দেয়। ADC সেটি digital মানে বদলায়; প্রসেসিং circuit মান গণনা করে; register সাময়িকভাবে রাখে; display circuit সংখ্যা দেখায়। ডিসপ্লেতে সংখ্যা দেখা গেলেই ভেতরের সেন্সর সরাসরি binary তথ্য দিচ্ছে—এ সিদ্ধান্ত ঠিক নয়।

